% LTeX: language=en-US
\section{IRFs}
\begin{frame}{The problem with IRFs}
\begin{witemize}
\item With a non-linear solution, traditional IRFs do not work anymore
\item The IRFs depend on two things:
\begin{enumerate}
  \item the sequence of future shocks, $u_t$ and
  \item the point in the state space at which the IRFs are started, i.e.\ the past history of shocks, $\Omega _t$.
\end{enumerate}
\item \textcite{KoopPesaranPotter1996} suggested the concept of Generalized IRFs (GIRF) that e.g.\ allow considering ``representative'' IRFs at the ergodic mean
\item The GIRF at time $t+n$ after a shock $u_t$ is given by
 \begin{equation}
 GIRF_n\left( {{u_t},{\Omega _{t - 1}}} \right) = E\left[ {{Y_{t + n}}|{u _t},{\Omega _{t - 1}}} \right] - E\left[ {{Y_{t + n}}|{\Omega _{t - 1}}} \right] \: , \label{eq:girf}
 \end{equation}
 (or some variation of this)
\item Given a point in the state space, the future shock realizations are averaged out
\end{witemize}
\end{frame}

\begin{frame}{How to generate GIRFs?}
\begin{witemize}
\item Additional problem: without pruning, IRF simulations can explode but with pruning one needs initial conditions not just for economic states, but for the whole augmented state space
\item There are infinitely many combinations of first to $n$th order terms that could have led to a particular value of the economic states
\item There are two special points one can consider:
\begin{enumerate}
  \item The ergodic mean where all components are at their expected values. i.e.\ the fixed point of the stochastic pruned solution
  \item The stochastic steady state, i.e.\ the fixed point of the deterministic pruned solution
\end{enumerate}

\end{witemize}
\end{frame}

\begin{frame}{How to generate GIRFs?}
\begin{witemize}
\item \textcite{AndreasenEtal2018} provide an analytical way to generate GIRFs from the pruned third order state space, starting at the ergodic mean
\item Dynare uses simulation-based approach to achieve the same result
\item The GIRFs can be computed from the difference between two simulated series:
\begin{enumerate}
  \item a simulated series of 140 periods using random shocks where at time $t=100$ an additional 1 standard deviation shock is added to whatever shock $\varepsilon_{100}$ was already drawn
  \item a simulated series of 140 periods using the same random shocks, but without the additional 1 standard deviation shock at time $t=100$
\end{enumerate}
\item By repeating this several times and taking the average, the effect of the random shocks is averaged out ($\approx$ Monte Carlo Integration)
\item For volatility shocks, a lot of repetitions may be necessary
\end{witemize}
\end{frame}

\begin{frame}{\textcite{FVetal2011AER} IRFs}
\begin{witemize}
\item \textcite{FVetal2011AER} do not compute GIRFs in the classical sense
\item They also condition on future shocks by setting them to 0 and starting the IRFs at the \highlight{ergodic mean in the absence of shocks/stochastic steady state}
\item Denote the future realization of shocks with $\Omega^{fut}$. The definition used is
 \begin{align}
  IR{F_n}\left( {{\nu _t},{\Omega _t}} \right) &= E\left[ {{Y_{t + n}}|{u_t},{\Omega _{t - 1}} = \left\{ { \ldots ,0} \right\},\Omega _{t + 1}^{fut} = \left\{ {0, \ldots } \right\}} \right] \notag \\
   &- E\left[ {{Y_{t + n}}|0,{\Omega _{t - 1}} = \left\{ { \ldots ,0} \right\},\Omega _{t + 1}^{fut} = \left\{ {0, \ldots } \right\}} \right]  \label{eq:FV_IRF} \: ,
\end{align}
where the expected values can be dropped as everything is deterministic
\item[$\rightarrow$] no repetitions necessary
\end{witemize}
\end{frame}

%\begin{frame}{The Problem With FV(2011) IRFs}
%\begin{itemize}
%\item Conditioning on future shocks and setting them to 0 only allows capturing part of the economic effects of risk shocks
%\item Because $\hat x_0^{f}=0$ and all higher order terms are based on this, the effect of the initial condition $ x - \bar x $ will mostly be neglected
%\item $\hat x_0^{f}=0$ follows from a first order policy function, which is known not to react to risk shocks
%\item Together with the conditioning on all shocks being 0 $\forall \: t+i, \: i>0$ this implies that only the ``risk adjustment channel'' is present via:
%\begin{itemize}
%  \item the constant term $1/2 \times g_{\sigma \sigma }\times \sigma ^2$
%  \item the time-varying risk-adjustment $1/2\times g_{u\sigma \sigma }\times  \sigma ^2 \times u_t$ in period $t$ where $u_t\neq0$ is present
%\end{itemize}
%\item \highlight{amplification effects} introduced by (risk) shocks and embedded in the other higher order terms are totally absent
%\item[$\rightarrow$] difference in the interaction between the location in the state space and future shocks is not captured
%\end{itemize}
%\end{frame}

%\begin{frame}{Dynare Output for Our RBC Model: SVD}
%\begin{figure}
%  \includegraphics[scale=.45]{../Codes/Identification/myRBCmodel_first_diff_identification/identification/myRBCmodel_first_diff_identification_ident_pattern_prior_mean_1}
%\end{figure}
%\end{frame}

\renewcommand{\Source}{}
